\documentclass[11pt]{article}
\usepackage[T1]{fontenc}
\usepackage{lmodern,microtype}
\usepackage[margin=.93in]{geometry}
\usepackage{amsmath,amssymb,amsthm,booktabs,hyperref}
\newtheorem{theorem}{Theorem}
\newtheorem{lemma}[theorem]{Lemma}
\newcommand{\Z}{\mathbb Z}
\newcommand{\Q}{\mathbb Q}
\setlength{\emergencystretch}{1.5em}
\title{Shifted binomial Hankel determinants, boundary degrees and Bernoulli coefficients}
\author{Blake Stone's mathematics project\\ \small Developed with OpenAI Codex}
\date{November 26, 2026}
\begin{document}\maketitle
\begin{abstract}
We determine the five determinant families in Cigler's Conjecture 17 on the range of nonnegative shift parameter $m$. Four of the identities continue to $m=-1$. The fifth gives a different polynomial at that boundary for every $k\ge1$. We give its exact finite determinant and prove that its degree is $3k-3$ for $k\ge2$, whereas the proposed expression has degree $k$. We also prove both degree predictions and the even Bernoulli leading coefficient in Conjecture 18, and give the corrected odd leading coefficient for all parameters. A Chebyshev moment determinant turns both evaluations into minors of Pascal's matrix.
\end{abstract}

\section{The identities}
For $r\ge0$ write $c_{r,j}=\binom{2j+r}{j}$ for $j\ge0$ and $c_{r,j}=0$ for $j<0$. Put
\[
 d_{r,s}(N)=\det(c_{r,i+j+s})_{0\le i,j<N},\qquad d_{r,s}(0)=1,\qquad
 \sigma_j=(-1)^{\binom j2}.
\]
Cigler and Tyson~\cite{ct} evaluated distinguished size classes in the unshifted case. Cigler~\cite[\S4, Conjecture 17]{cigler} proposed five identities with $s=-m$ for $m\ge-1$.
\begin{theorem}\label{main}
Let $k\ge1$, $m\ge-1$, $K=k+m$, and $n\ge0$. With $r=2K+1$,
\begin{align}
 d_{r,-m}(rn)&=(2n+1)^k,\label{odd0}\\
 d_{r,-m}(rn+m+1)&=\sigma_{m+1}(2n+1)^k.\label{odd1}
\end{align}
If $m\ge0$, then also
\begin{equation}\label{third}
 d_{r,-m}(rn+K+1)=\sigma_{K+1}4^k(n+1)^k.
\end{equation}
With $r=2K$,
\begin{align}
 d_{r,-m}(rn)&=(-1)^{Kn},\label{even0}\\
 d_{r,-m}(rn+m+1)&=\sigma_{m+1}(-1)^{Kn}.\label{even1}
\end{align}
At $m=-1$ the left side of \eqref{third} is $2n+3$ for $k=1$, and it is a polynomial of degree $3k-3$ in $n$ for $k\ge2$. Formula~\eqref{boundary} below evaluates it for every $k,n$.
\end{theorem}
The three odd identities in their stated ranges also follow directly from Chern and Shi~\cite[Theorem 1.9 and (7.1)]{cs}. Set their height to $K$, their shift to $k$ and $c=2$. Their entry $\operatorname{Cat}^{(1,2)}_{u,K}=\binom{2u+1}{u+K+1}$, with $u=i+j+k$, becomes $c_{2K+1,i+j-m}$ by binomial symmetry. The first two size classes require $K\ge k-1$ and the third requires $K\ge k$. These are respectively $m\ge-1$ and $m\ge0$. The proof below is independent of that specialization and determines the exceptional class excluded by its hypothesis. Their equation (1.12) also gives the particular $k=1$ boundary value $2n+3$.

For example, $k=2,m=-1,n=0$ in \eqref{third} asks for the size-two determinant of the entries $5,21,84$. It is $5\cdot84-21^2=-21$. The proposed value is $-16$. The boundary polynomial at $k=2$ is $-(n+1)(2n+3)(6n+7)$.

\section{An endpoint and root determinant}
Take $r>0$; the only case $r=0$ in the theorem is $k=1,m=-1$, and both even identities then read $d_{0,1}(0)=1$. With $y=2\cos\theta$ and $\theta$ uniform on $[0,\pi]$, let $L(f)$ be the mean of $f(y^2)$. Write $T_j(\cos\theta)=\cos(j\theta)$. Orthogonality of Fourier frequencies gives, for $2j+r\ge0$,
\[
 \frac1\pi\int_0^\pi(2\cos\theta)^{2j+r}\cos(r\theta)\,d\theta=c_{r,j}.
\]
Indeed the coefficients of the two frequencies $\pm r$ in $(e^{i\theta}+e^{-i\theta})^{2j+r}$ both equal $\binom{2j+r}{j}$, or both vanish if $j<0$. If $r$ is even set $a=k$; if $r$ is odd set $a=k+1$. We have
\[
 2y^{r-2m}T_r(y/2)=M(y^2),\quad
 M(x)=x^a\prod_{i=1}^K(x-\alpha_i),\quad
 \alpha_i=4\cos^2\!\left(\frac{(2i-1)\pi}{2r}\right).
\]
This follows from the leading coefficient and roots of $2T_r(y/2)$. Put $d=a+K$. The inputs $c_{r,j-m}$ are the moments $L(x^j M(x))/2$ for $j\ge0$.

Let $Q_j(x)=2T_{2j}(\sqrt{x}/2)$ for $j\ge0$. These are monic for $j\ge1$, with norm squared two for $L$. The monic constant polynomial is $1=Q_0/2$ with norm squared one. Define $E(N)$ from columns $Q_N,\ldots,Q_{N+d-1}$ by taking coefficient rows $[x^v]$ for $0\le v<a$ and evaluations at the $\alpha_i$.
\begin{lemma}\label{chris}
For $N\ge0$, \(d_{r,-m}(N)=(-1)^{Nd}\det E(N)/\det E(0)\).
\end{lemma}
\begin{proof}
This is a specialization of the Christoffel determinant formula; compare \cite{kratt}. Use the ordered rows $L(x^i\,\cdot)$ for $i<N$ followed by the $d$ endpoint and root functionals. In the monic orthogonal basis of degrees $0,\ldots,N+d-1$, the first $N$ rows vanish on the last $d$ columns and have triangular first block with diagonal given by the squared norms. In the ordered monic basis
\[
 M,xM,\ldots,x^{N-1}M,\quad1,x,\ldots,x^{d-1}
\]
the bottom left block vanishes, the top left block has determinant $2^N d_{r,-m}(N)$, and the bottom right block is the evaluation matrix for degrees below $d$. These degrees have been reordered by sign $(-1)^{Nd}$. The bottom right determinant is nonzero since the common kernel of the $d$ functionals consists of multiples of $M$. For $N\ge1$ the product of the squared norms is $2^{N-1}$. Replacing the monic constant by $Q_0=2$ doubles the reference determinant at zero, canceling the remaining factor $1/2$. This gives the formula, also valid by definition at $N=0$.
\end{proof}

Write $\phi_i=(2i-1)\pi/r$. Replace $E(N)$ by a row equivalent matrix whose column of frequency $v=N+j$ has entries
\[
 (-1)^v v^{2h}\ (0\le h<a),\qquad \cos(v\phi_i)\ (1\le i\le K),
 \qquad 0^0=1.
\]
The row transformation is independent of $N$. To see this, set $x=2-2\cosh u$. Then $Q_v(x)=2(-1)^v\cosh(vu)$, by the cosine identity and polynomial continuation. Extracting $[u^{2h}]$ in place of $[x^h]$ is triangular and invertible since $x=-u^2+O(u^4)$. At a positive root, $Q_v(\alpha_i)=2\cos(v\phi_i)$. Use the resulting ratio in Lemma~\ref{chris} also at negative $N$. Since $Q_{-v}=Q_v$, reversing the columns gives
\begin{equation}\label{reflect}
 d_{r,-m}(1-d-N)=\sigma_d d_{r,-m}(N)\qquad(N\in\Z)
\end{equation}
for this extension. The prefactor change has sign $(-1)^{d(1-d)}=1$.

\section{Pairing frequencies}
First let $r=2K$. At $N=rn$, the columns with offsets $0,\ldots,K-1$ form an invertible root block up to row factors $(-1)^n$, since their initial polynomials have consecutive degrees. Offset $K$ vanishes on the roots. For $1\le z<k$, add column $K-z$ to column $K+z$. This cancels the root entries since $\cos((r-j)\phi_i)=-\cos(j\phi_i)$; the partner exists even when $m=-1$. For $M_0=N+K$ the endpoint entries in these $k$ columns, apart from their individual signs $(-1)^{N+K+z}$, are
\[
 (M_0+z)^{2h}+(M_0-z)^{2h},\quad 0\le h<k,\ 0\le z<k,
\]
with the $z=0$ column halved. As polynomials in $z^2$ these rows have degrees $h$ with leading coefficient two. The resulting Vandermonde determinant is independent of $N$. The root row factors give $(-1)^{Kn}$, and $N$ is even, so the remaining sign changes from $N=0$ are all one. This proves \eqref{even0} for every integer $n$.

Now let $r=2K+1$. Group row $h=0$ with the $K$ root rows. Their initial square block on offsets $0,\ldots,K$ is invertible by evaluation at the $K+1$ distinct nodes $0,\alpha_1,\ldots,\alpha_K$. For $1\le j\le k$ set
\[
 S_h(A,z)=(A+z)^{2h}-(A-z)^{2h},\qquad
 V(A;z_1,\ldots,z_k)=\det(S_h(A,z_j))_{1\le h,j\le k}.
\]
After extracting $z_j$ in each column the entries become polynomials in $z_j^2$ of degree $h-1$ with leading coefficient $4hA$. Hence
\begin{equation}\label{v}
 V(A;\mathbf z)=(4A)^k k!\left(\prod_j z_j\right)
                       \prod_{i<j}(z_j^2-z_i^2).
\end{equation}

At $N=rn$, add the column at offset $K+1-j$ to that at $K+j$. Their values on the first group cancel, since the offsets sum to $r$ and have opposite parity. Including $m=-1$ the partners range from $K$ down to $K+1-k\ge0$. The last $k$ columns on rows $h=1,\ldots,k$, with signs $(-1)^{N+K+j}$ taken out, give $V(N+r/2;1/2,3/2,\ldots,k-1/2)$. Relative to $N=0$, the first group contributes $(-1)^{n(K+1)}$, the extracted signs $(-1)^{Nk}$, and the prefactor $(-1)^{Nd}$. Their product is one since $d=K+k+1$ and $N\equiv n\pmod2$. Formula~\eqref{v} therefore gives \eqref{odd0} for every integer $n$.

In either parity $r-d+1=m+1$. Reflect the size $rn+m+1$ in \eqref{reflect} to obtain $-r(n+1)$. Substitution of \eqref{odd0} or \eqref{even0} gives \eqref{odd1} or \eqref{even1}. The sign reductions use $\binom{2k+t}{2}\equiv\binom t2+k\pmod2$.

For the third size take $N=rn+K+1$ and $A=r(n+1)=N+K$. In the first $K+1$ columns, the grouped rows are $(-1)^{n+1}$ times their initial values with offsets $K,K-1,\ldots,0$. For $m\ge0$ every column at $K+j$, $1\le j\le k$, has a partner at $K-j$. Subtract the partner to clear the grouped rows. The remaining block, after removing column signs, is $V(A;1,\ldots,k)$. Relative to $N=0$, the full sign is
\[
 (-1)^{N(d+k)+(n+1)(K+1)}\sigma_{K+1}=\sigma_{K+1}.
\]
The final factor is the reversal of $K+1$ columns in the grouped rows. Thus the required determinant is $\sigma_{K+1}V(A;1,\ldots,k)/V(r/2;1/2,\ldots,k-1/2)$. In \eqref{v}, the ratio of the products in the $z$ variables for the integer and half-integer lists is $2^k$. Indeed on appending the $j$th element this ratio gains
\[
 \frac{j}{j-\frac12}\prod_{i=1}^{j-1}\frac{j+i}{j+i-1}=2.
\]
The ratio of the $A$-powers supplies another $2^k(n+1)^k$, proving \eqref{third}.

\section{The exceptional family}
If $m=-1$ then $K=k-1$, so the partner at offset $K-k$ is absent. The last offset equals $r$ and has opposite grouped entries to offset zero. Adding column zero to the last column clears these entries. Keep the same extracted column sign as before. With $A=r(n+1)$, the last entry in row $h$ then equals
\[
 (N+r)^{2h}-N^{2h}
   =S_h(A,k)-S_h(A-r/2,1/2).
\]
Indeed $A+k=N+r$ and $A-k=N-1$. Write $V_*(A)$ for the determinant with its first $k-1$ columns $S_h(A,j)$ and last column as displayed. All previous signs and reference factors still apply. Hence
\begin{equation}\label{boundary}
 d_{2k-1,1}((2k-1)n+k)=
 \sigma_k\,\frac{V_*\big((2k-1)(n+1)\big)}
 {V((2k-1)/2;1/2,3/2,\ldots,k-1/2)} .
\end{equation}
At $k=1$ this is $2n+3$. At $k\ge2$ we determine its exact degree. Expand each column in the vectors
\[
 \left(\binom{2h}{t}A^{2h-t}\right)_{h=1}^k,\quad 1\le t\le2k,
\]
regarding entries with $t>2h$ as zero. Ordinary columns $j<k$ have coefficients $2j^t$ for odd $t$, and zero for even $t$. The exceptional column has coefficient $r\ne0$ at $t=2$ by the binomial expansion of $-S_h(A-r/2,1/2)$. A nonzero determinant needs distinct $t$'s. The minimum possible sum is $(k-1)^2+2$, attained only by the $k-1$ odd indices $1,3,\ldots,2k-3$ together with $2$. The coefficient minor on the column side for that set is nonzero by the Vandermonde determinant on $1^2,\ldots,(k-1)^2$.

The minor $\det(\binom{2h}{t})$ for these sorted indices is positive. For completeness, $\binom{2h}{t}$ counts paths from $(0,t)$ to $(2h,0)$ with steps $(1,0)$ and $(1,-1)$. In the determinant expansion, intersecting path families cancel by the usual tail swap at their first intersection. A disjoint family must pair the sorted indices in order, since reversed endpoints force an intersection. For these diagonal pairings $t_i\le 2i$. Take each path horizontally first and then down; at a common abscissa $x$ its height is $\min(t_i,2i-x)$, strictly increasing with $i$. This gives a disjoint family and proves positivity.

The leading power of $A$ is consequently $\sum_{h=1}^k2h-(k-1)^2-2=3k-3$. Since the denominator in \eqref{boundary} is a nonzero constant and $A$ is linear in $n$, this proves Theorem~\ref{main}.
\section{The two Bernoulli leading coefficients}
Let $B_j$ be the Bernoulli numbers defined by $\sum B_jz^j/j!=z/(e^z-1)$. For $k\ge2,m\ge0$ put $K=k+m$. Cigler's next conjecture concerns
\[
 L_{\rm e}(n)=(-1)^{Kn}d_{2K,-m}(2Kn+m+2),\qquad
 L_{\rm o}(n)=d_{2K+1,-m}((2K+1)n+m+2).
\]
\begin{theorem}\label{bernoulli}
These are polynomials in $n$, of respective degrees $2k-3$ and $3k-3$. With $\tau=-\sigma_m$ their leading coefficients are
\begin{align*}
 [n^{2k-3}]L_{\rm e}
 &=\tau\,\frac{4^{2k-2}K^{2k-3}(2^{2k-2}-1)B_{2k-2}}{(2k-2)!},\\
 [n^{3k-3}]L_{\rm o}
 &=\tau\,\frac{2^{3k-2}(2K+1)^{2k-2}B_{2k-2}}{(2k-2)!}.
\end{align*}
\end{theorem}
The first is the even formula in Conjecture 18 of \cite{cigler}. The odd formula there has $2k+m+1$ in place of $2K+1=2k+2m+1$. For instance, at $k=2,m=1$ the original odd determinants for $n=0,1,2,3$ are $-25,-346,-1355,-3444$. Their third difference $-392$ gives leading coefficient $-196/3$, whereas the printed formula gives $-48$.

\begin{proof}
In either parity reflect the size $rn+m+2$ by \eqref{reflect}. Its partner is $rt-1$, where $t=-n-1$. Write $A=rt+r/2=-r(n+1/2)$. The columns have frequencies $rt+j$ for offsets $-1,0,\ldots,d-2$. Move offset $-1$ to the end and subtract the column at $+1$ from it. This clears its root entries since the root values at $rt\pm1$ agree, and clears its row $h=0$ too. The remaining exceptional column after extracting $(-1)^{rt}$ is
\[
 S_h(rt,1)=S_h(A-r/2,1).
\]
The ordinary columns are the earlier $N=rt$ calculation with the last paired column removed. The sign from moving the exceptional column is $(-1)^{d-1}$, precisely the sign belonging to the removed column in the $N=0$ reference matrix (its offset is $d-1$). Compared with size $rt$, the external prefactor has gained $(-1)^d$.

For odd $r=2K+1$, let $W_j=\prod_{u=1}^j z_u\prod_{u<v}(z_v^2-z_u^2)$ for $z_u=u-1/2$, where the products use $1\le u<v\le j$. Let $V_{\rm o}(A)$ have ordinary columns $S_h(A,z_u)$, $1\le u<k$, followed by the exceptional column, with $1\le h\le k$. The preceding signs and \eqref{reflect} give
\[
 L_{\rm o}(n)=\sigma_d(-1)^d
 \frac{\det V_{\rm o}(A)}{4^k k!(r/2)^kW_k}.
\]
The expansion in vectors $\binom{2h}{j}A^{2h-j}$ from the preceding section applies again. The least sum of indices uses $1,3,\ldots,2k-3$ for the ordinary columns and $2$ for the exceptional one. Its coefficient at $j=2$ is
$(-r/2+1)^2-(-r/2-1)^2=-2r$.
Order the selected indices with $2$ last. Their row minor is the determinant obtained from the lower triangular matrix $(\binom{2h}{2v-1})_{h,v=1}^k$ by replacing its last column with $\binom{2h}{2}$. Write $c_j$ for the coefficients of the expansion
\[
 \binom{2h}{2}=\sum_{j=1}^h c_j\binom{2h}{2j-1}.
\]
Exponential generating functions give
\[
 \sum_{j\ge1}c_j\frac{x^{2j-1}}{(2j-1)!}
 =\frac{x^2}2\coth x,\qquad
 c_k=(2k-1)2^{2k-3}B_{2k-2}.
\]
Thus the replaced row determinant is $2^k k!c_k$. The column minor is $(-2r)2^{k-1}W_{k-1}$. Their product is the coefficient of $A^{3k-3}$ in $\det V_{\rm o}(A)$, and higher coefficients vanish by the minimum sum of indices. Since $W_k/W_{k-1}=(2k-1)!/2$, the conversion $A=-r(n+1/2)$ gives the stated odd coefficient. Its sign outside $B_{2k-2}$ follows from $\sigma_d(-1)^{d+3k-2}=-\sigma_m$ for $d=2k+m+1$.

For the even case $r=2K$, put $Z_j=\prod_{0\le u<v<j}(v^2-u^2)$. Let $V_{\rm e}(A)$ have ordinary columns
\[
 (A+z)^{2h}+(A-z)^{2h},\qquad z=0,\ldots,k-2,
\]
with the zero-offset column halved, followed by $S_h(A-r/2,1)$; here $0\le h<k$. The even residue evaluation contributes $(-1)^{Kt}$, which combines with the defining factor $(-1)^{Kn}$ to give $(-1)^K$. Hence
\[
 L_{\rm e}(n)=\sigma_d(-1)^{d+K}
             \frac{\det V_{\rm e}(A)}{2^{k-1} Z_k}.
\]
The least sum of distinct indices now uses $0,2,\ldots,2k-4$ for the ordinary columns and $1$ for the exceptional one, whose coefficient at $j=1$ is $2$. This gives degree at most $2k-3$. Order these indices with $1$ last. In the unit lower triangular basis $(\binom{2h}{2j})_{h,j=0}^{k-1}$ the last-column replacement $\binom{2h}{1}$ gives determinant $b_{k-1}$, where
\[
 \sum_{j=0}^h b_j\binom{2h}{2j}=\binom{2h}{1},\qquad
 \sum_{j\ge0} b_j\frac{x^{2j}}{(2j)!}=x\tanh x,\qquad
 b_{k-1}=2^{2k-2}(2^{2k-2}-1)B_{2k-2}.
\]
The column minor is $2^{k-1}Z_{k-1}$. Now $Z_k/Z_{k-1}=(2k-2)!/2$ and $\sigma_d(-1)^{d+K+1}=-\sigma_m$ for $d=2k+m$. Substituting $A$ proves the even formula. Both leading coefficients are nonzero because $B_{2j}\ne0$ for $j\ge1$.
\end{proof}

\begin{thebibliography}{9}
\bibitem{cigler} J. Cigler, \emph{Some experimental observations about Hankel determinants of convolution powers of Catalan numbers}, v2, \S4, Conjecture 17 (2023). \href{https://arxiv.org/abs/2308.07642v2}{arXiv:2308.07642v2}.
\bibitem{ct} J. Cigler and M. Tyson, \emph{An interesting class of Hankel determinants}, in Algorithmic Combinatorics (2020), 137--157. \href{https://arxiv.org/abs/1807.08330}{arXiv:1807.08330}.
\bibitem{cs} S. Chern and W. Shi, \emph{Hankel determinants of Catalan-like sequences} (2026). \href{https://arxiv.org/abs/2608.27208}{arXiv:2608.27208}.
\bibitem{kratt} C. Krattenthaler, \emph{Hankel determinants of linear combinations of moments of orthogonal polynomials, II}, Ramanujan J. \href{https://doi.org/10.1007/s11139-021-00514-8}{doi:10.1007/s11139-021-00514-8}.
\end{thebibliography}
\end{document}
